% capitulos/03-recursos-financeiros.tex

Esta pesquisa conta com financiamento de fontes públicas de fomento
à ciência e tecnologia, conforme detalhado nas seções a seguir. O
gerenciamento dos recursos é realizado pela Fundação de Empreendimentos
Científicos e Tecnológicos (\abreviatura{FINATEC}{Fundação de
Empreendimentos Científicos e Tecnológicos}), instituição gestora
vinculada à Universidade de Brasília.

\section{Fontes de Financiamento}

O projeto é financiado por meio de duas fontes principais:

\begin{enumerate}
  \item \textbf{Bolsa de Doutorado CAPES} (Processo n.º 0000000/2023-0):
        bolsa de pesquisa no valor de R\$~2.200,00/mês, com vigência de
        48 meses (março/2023 a fevereiro/2027);
  \item \textbf{Edital Universal CNPq} (Processo n.º 000000/2023-0):
        auxílio à pesquisa no valor de R\$~50.000,00, com vigência de
        24 meses.
\end{enumerate}

\section{Orçamento Detalhado}

A \autoref{tab:orcamento} apresenta o orçamento previsto para a execução
da pesquisa, discriminado por categoria de despesa.

\begin{table}[H]
  \centering
  \caption{Orçamento previsto para execução da pesquisa
           (valores em R\$, referência: janeiro/2025).}
  \label{tab:orcamento}
  \begin{tabularx}{\textwidth}{Xrrr}
    \toprule
    \textbf{Item / Descrição}
      & \textbf{Qtd.}
      & \textbf{Unit. (R\$)}
      & \textbf{Total (R\$)} \\
    \midrule
    \multicolumn{4}{l}{\textit{Material de Consumo}} \\
    \thinrule
    \quad Frascos de polietileno (250 mL)      & 500 & 3,50   & 1.750,00 \\
    \thinrule
    \quad Sacos plásticos para coleta (1 L)    & 300 & 1,20   & 360,00   \\
    \thinrule
    \quad Reagentes laboratoriais (kit)        &   3 & 800,00 & 2.400,00 \\
    \thinrule
    \quad Material de escritório e impressão   &   1 & 500,00 & 500,00   \\
    \midrule
    \multicolumn{4}{l}{\textit{Serviços de Terceiros}} \\
    \thinrule
    \quad Análises laboratoriais externas      & 270 & 45,00  & 12.150,00 \\
    \thinrule
    \quad Cartografia / SIG (consultoria)      &  40 & 80,00  & 3.200,00 \\
    \midrule
    \multicolumn{4}{l}{\textit{Diárias e Passagens}} \\
    \thinrule
    \quad Diárias (trabalho de campo, 3 campanhas) & 15 & 340,00 & 5.100,00 \\
    \thinrule
    \quad Passagens aéreas (eventos científicos)   &  4 & 800,00 & 3.200,00 \\
    \midrule
    \multicolumn{4}{l}{\textit{Equipamentos}} \\
    \thinrule
    \quad GPS de precisão (receptor GNSS)      &   1 & 4.500,00 & 4.500,00 \\
    \thinrule
    \quad Coletor de solos tipo Uhland         &   2 & 650,00   & 1.300,00 \\
    \midrule
    \multicolumn{4}{l}{\textit{Publicações}} \\
    \thinrule
    \quad Taxas de publicação (APC) em periódicos  &  3 & 4.000,00 & 12.000,00 \\
    \thinrule
    \quad Revisão e tradução de manuscritos        &  3 & 1.500,00 &  4.500,00 \\
    \midrule
    \textbf{TOTAL GERAL}                       & & & \textbf{50.960,00} \\
    \bottomrule
  \end{tabularx}
  \fonte{Elaborado pelo autor (2025).}
\end{table}

\section{Contrapartida Institucional}

A Universidade de Brasília contribui com infraestrutura laboratorial,
espaço físico, acesso a bases de dados científicas e supervisão acadêmica
sem ônus ao projeto. O valor estimado da contrapartida institucional é
de R\$~15.000,00, correspondendo a horas de uso de equipamentos
multiusuários do \abreviatura{LQAA}{Laboratório de Química Ambiental
e Agroecologia}.
